\chapter{Typical Use}
\section{Executable}
In an executable, just declare a static constexpr
\texttt{Dwm::What::Info} in the same translation unit as
\texttt{main()}, preferably in a unique namespace.  You may want to
mark your instance with \texttt{\_\_attribute\_\_((used))} to prevent
the toolchain from optimizing it out of existence due to lack of use.
For example:

\begin{verbatim}
namespace MyNamespace {
   namespace myproginfo {
      static constexpr __attribute((used))
          info(DWM_WHAT_TYPE_EXE, DWM_WHAT_STATUS_REL, "myprog", "1.0.0",
               @DWM_WHAT_SYM_COPYRIGHT "My Name 2025", "my program");
   }
}
\end{verbatim}

\section{Compiled library}
In one of your source files, just declare a static constexpr
\texttt{Dwm::What::Info}, preferably in a unique namespace.  You may
want to mark your instance with \texttt{\_\_attribute\_\_((used))}
to prevent the toolchain from optimizing it out of existence due to
lack of use.  For example:

\begin{verbatim}
namespace MyNamespace {
   namespace mylibinfo {
      static constexpr  __attribute__((used))
         info(DWM_WHAT_TYPE_LIB, DWM_WHAT_STATUS_REL, "mylib", "1.0.0",
              @DWM_WHAT_SYM_COPYRIGHT "My Name 2025", "my library");
   }
}
\end{verbatim}

\section{Header-only library}
In one of your header files, declare an inline constexpr
\texttt{Dwm::What::Info}, preferably in a unique namespace.  You may
want to mark your instance with \texttt{\_\_attribute\_\_((used))}
to prevent the toolchain from optimizing it out of existence due to
lack of use.  For example:

\begin{verbatim}
namespace MyNamespace {
   namespace myhdrinfo {
      inline constexpr  __attribute__((used))
         info(DWM_WHAT_TYPE_HDR, DWM_WHAT_STATUS_REL, "myheaderlib", "1.0.0",
              @DWM_WHAT_SYM_COPYRIGHT "My Name 2025", "my header library");
   }
}
\end{verbatim}

\subsection{Caveats}
It's worth noting that the usual caveats apply here.  Despite the fact that
this is the recommended means of producing a single instance in a linked
program, ODR (One Definition Rule) violations can rear their head in
difficult-to-diagnose ways.  Your header is making a promise it can't
really keep (unless the header never changes), since each compilation
unit is going to get its own copy and the linker is going to pick one.
If you happen to have an object file in a old library that was compiled
against an older version, and link it to an object file that was compiled
against a newer version, the linker is going to just pick one, with
NDR (No Diagnostic Required) for ODR violations.  My recommendation is
to name your instance differently for each release, and provide an
accessor that returns a reference to that instance.  Then the linker
will see two identifiers instead of one, and will keep both.

I mention this because in the case of
\texttt{Dwm::What::SegmentedLiteral} and \texttt{Dwm::What::Info}, a
declaration is typically not just stamping out a variable, but a new
type.  The template is parameterized by the number and size of the
segments, so any change to either of these will yield a different
type, not just different content.  The linkers of today don't seem to
use this information to flag ODR violations.
